Fix an arbitrary prescribed \(\chi>1\). Set \(K=4\), \(c_1=c_2=(5\chi-3)/4\), \(c_3=c_4=1/2\), \(b_E=1\), \(b_O=2\), and \(\kappa=(5\chi-3)/2\), and put
\[
 y_{\chi}=\sqrt{1-\frac{3}{10\chi}}.
\]
Let \(X\) and every \(S_{2d}\) be constant. Set \(P_{\epsilon}(G=E)=P_{\epsilon}(G=O)=1/2\), and let a Rademacher \(U\) be independent of \(G\). In source \(E\), \(P_{\epsilon}(T=t\mid U=u,G=E)=1/2\); in source \(O\), \(P_{\epsilon}(T=t\mid U=u,G=O)=(1+tu/4)/2\). The proxies \(W_{\epsilon},W_{y_{\chi}},W_{1/2}^{(1)},W_{1/2}^{(2)}\) are mutually independent conditional on \(U\), are independent of \(G,A\) conditional on \(U\), and form
\[
 Z(a)=(W_{\epsilon},W_{y_{\chi}},W_{1/2}^{(1)},W_{1/2}^{(2)}),
\]
while \(Y(a)=U\) for both treatments. Take \(\mathcal D=\{d_{\mathrm p},d_{\mathrm e}\}\), where \(d_{\mathrm p}\) assigns \(W_{\epsilon}\) to role one and \(W_{y_{\chi}}\) to role three, whereas \(d_{\mathrm e}\) assigns \(W_{1/2}^{(1)}\) to role one and \(W_{1/2}^{(2)}\) to role three. Then
\[
 c_{E,d_{\mathrm p}}=\frac{5\chi-1}{2},\qquad
 c_{O,d_{\mathrm p}}=\frac{5\chi+1}{2},\qquad
 c_{E,d_{\mathrm e}}=2,\qquad c_{O,d_{\mathrm e}}=3,
\]
so \(C_{d_{\mathrm p}}=5\chi\) and \(C_{d_{\mathrm e}}=5\). For a design with early signal \(x\) and late signal \(y\), the candidate bridges are
\[
 h_d(l,a)=\frac{l}{y},
 \qquad
 q_d(e,a)=\frac{1-(2a-1)e/(4x)}{1-1/16}.
\]
For later center calculations, name the scalar
\[
 R(x,y):=\frac{(1-x^2)(1-y^2)}{1-x^2y^2}.                 \tag{1}
\]
Lemma~\(\mathrm{lem{:}sharp\text{-}uniform\text{-}binary\text{-}elbow}\), rather than this construction, proves that the unrestricted unadjusted risk and the distinct CLW affine linear residual risk both equal \(R(x,y)\) at the symmetric centers.  No such equality is imposed on perturbations in the constrained model.  The named transported CLW selector compares \(C_dr_d^{\mathrm{lin}}(P)/B\).  The generic unrestricted-risk selectors remain separate secondary comparisons.

The \emph{binary constrained finite-cell model} \(\mathcal M_{\mathrm{bin}}(\chi)\) is the family of laws \(P\) with the preceding prescribed catalogue, costs, cap, and role assignments, constant \(X\) and constant \(S_{2d}\), latent variables \(U,T,W_1,W_2,W_3,W_4\in\{-1,1\}\), \(A=(T+1)/2\), and potential variables
\[
 Z(a)=(W_1,W_2,W_3,W_4),\qquad Y(a)=U,
 \qquad a\in\{0,1\}.
\]
For strictly positive probability vectors \((\pi_g:g\in\{E,O\})\), \((\nu_u:u\in\{-1,1\})\), \((\lambda_t:t\in\{-1,1\})\), observational treatment kernels \(e_u(t)\), and proxy kernels \(k_{j,u}(w)\), its latent-cell mass factorizes as
\[
 P(G=g,U=u,T=t,W_1=w_1,\ldots,W_4=w_4)
 =\pi_g\nu_u\bigl\{\mathbf 1\{g=E\}\lambda_t
 +\mathbf 1\{g=O\}e_u(t)\bigr\}
 \prod_{j=1}^4k_{j,u}(w_j).                               \tag{2}
\]
Thus \(U\) is independent of \(G\), treatment is randomized independently of \(U\) in source \(E\), and the four proxies are mutually independent conditional on \(U\) and invariant across source and treatment. The design \(d_{\mathrm p}\) assigns \((W_1,W_2)\) to roles one and three, and \(d_{\mathrm e}\) assigns \((W_3,W_4)\) to roles one and three; role two is constant. For each \(0<\epsilon<10^{-3}\), the center \(P_\epsilon\) is the point having \(\pi_E=\pi_O=1/2\), \(\nu_{-1}=\nu_1=1/2\), \(\lambda_{-1}=\lambda_1=1/2\), \(e_u(t)=(1+tu/4)/2\), and
\[
 k_{j,u}(w)=\frac{1+s_juw}{2},
 \qquad
 (s_1,s_2,s_3,s_4)=(\epsilon,y_{\chi},1/2,1/2).
\]
The model is parameterized by the relative interiors of these finite probability simplices. Its relative topology is equivalently the Euclidean parameter topology or the maximum-norm topology \(\|\cdot\|_{\mathrm{cell}}\) induced by the complete latent-cell vector in (2).